\chapter{The command line}
\label{ch:shell}

A \textbf{shell} reads what you type and runs it. On Windows you will probably use PowerShell to connect; on the cluster, you will be using \texttt{bash}. macOS and Linux terminals work almost the same as \texttt{bash}. For a longer course, see Shotts (2019) or \href{https://swcarpentry.github.io/shell-novice/}{Software Carpentry}\footnote{\texttt{https://swcarpentry.github.io/shell-novice/}}.

\section{Anatomy of a command}\label{anatomy-of-a-command}

\begin{lstlisting}[language=bash]
# cluster
ls -l ~/project
\end{lstlisting}

\texttt{ls} is the command: what to do (list files). \texttt{-l} is an option: how to do it (long format). \texttt{\textasciitilde{}/project} is the argument: what to act on. Spaces separate the parts, so a path containing a space needs double quotes, or the shell reads it as two things.

\section{Paths}\label{paths}

{\def\LTcaptype{none} % do not increment counter
\begin{longtable}[]{@{}ll@{}}
\toprule\noalign{}
Symbol & Meaning \\
\midrule\noalign{}
\endhead
\bottomrule\noalign{}
\endlastfoot
\texttt{\textasciitilde{}} & Your home folder \\
\texttt{.} and \texttt{..} & This folder, and the one above \\
\texttt{/} & Top of the Linux file tree, and the separator between folders \\
\texttt{*} & Wildcard: \texttt{case\_*.json} matches every such file \\
\texttt{\$HOME}, \texttt{\$USER} & Variables that hold your home folder and username \\
\end{longtable}
}

Linux names are case-sensitive, so \texttt{Data} and \texttt{data} differ. Windows writes \texttt{C:\textbackslash{}Users\textbackslash{}me}, and Linux writes \texttt{/users/me}.

\section{Commands worth knowing}\label{commands-worth-knowing}

{\def\LTcaptype{none} % do not increment counter
\begin{longtable}[]{@{}
  >{\raggedright\arraybackslash}p{(\linewidth - 2\tabcolsep) * \real{0.5000}}
  >{\raggedright\arraybackslash}p{(\linewidth - 2\tabcolsep) * \real{0.5000}}@{}}
\toprule\noalign{}
\begin{minipage}[b]{\linewidth}\raggedright
Command
\end{minipage} & \begin{minipage}[b]{\linewidth}\raggedright
What it does
\end{minipage} \\
\midrule\noalign{}
\endhead
\bottomrule\noalign{}
\endlastfoot
\texttt{pwd} & Print the folder you are in \\
\texttt{ls}, \texttt{ls\ -l}, \texttt{ls\ -t} & List files; with details; newest first \\
\texttt{cd\ folder}, \texttt{cd\ ..}, \texttt{cd} & Go to a folder; up one; home \\
\texttt{mkdir\ -p\ a/b} & Make folders, including parents, with no error if they exist \\
\texttt{cp\ a\ b}, \texttt{cp\ -r\ a\ b} & Copy a file; \texttt{-r} copies a folder and its contents \\
\texttt{mv\ a\ b} & Move or rename \\
\texttt{rm\ file} & Delete. \textbf{There is no recycle bin.} Look first with \texttt{ls} \\
\texttt{ln\ -s\ target\ name} & Make a shortcut (a symbolic link) \\
\texttt{cat}, \texttt{head\ -5}, \texttt{tail\ -20} & Print a whole file; its first 5 lines; its last 20 \\
\texttt{less\ file} & Read a long file page by page; \texttt{q} quits \\
\texttt{wc\ -l} & Count lines \\
\texttt{grep\ "text"\ file} & Show lines containing text. \texttt{-i} ignores case, \texttt{-l} lists only file names \\
\texttt{sed\ -i\ \textquotesingle{}s/old/new/\textquotesingle{}\ file} & Replace text in a file \\
\texttt{chmod\ u+x\ file} & Let yourself (\texttt{u}) execute (\texttt{x}) a file. \texttt{chmod\ -R\ u+w\ folder} lets you write to everything in a folder \\
\texttt{tar\ czf\ a.tgz\ folder} & Pack a folder into one compressed archive (create, gzip, file) \\
\texttt{tar\ xzf\ a.tgz} & Unpack it (extract) \\
\texttt{df\ -h\ .}, \texttt{du\ -sh\ folder} & Free space on this disk; size of a folder \\
\texttt{hostname}, \texttt{nproc} & Which computer am I on? How many cores may I use? \\
\texttt{man\ ls} & The manual for a command; \texttt{q} quits \\
\end{longtable}
}

To edit on the cluster, \texttt{nano\ file} is the easiest editor if it is installed: the keys are listed at the bottom. NI-HPC\textquotesingle s guide uses \texttt{vi}, which is powerful but odd at first.

\section{Joining commands}\label{joining-commands}

\begin{lstlisting}[language=bash]
# cluster

# | (a pipe) feeds one command's output into the next: count the log files
ls logs/*.out | wc -l

# > saves output to a file, replacing what was there; >> adds to the end
ls -l > listing.txt

# && runs the second command only if the first succeeded
mkdir -p results && cd results

# ; runs the second command whatever happened to the first
cd results; ls

# 2>/dev/null throws away error messages (stream 2) instead of printing them
ls /no/such/folder 2>/dev/null

# $( ) puts one command's output inside another: show the newest log
cat "$(ls -t logs/*.out | head -1)"
\end{lstlisting}

\section{Variables and loops}\label{variables-and-loops}

\begin{lstlisting}[language=bash]
# cluster

# Set a variable. No spaces around the =
DATA=/mnt/scratch2/users/$USER/data

# Use it with $. Double quotes keep it in one piece even if it contains a space
ls "$DATA"

# A for loop: repeat for each folder; $d holds the current one
for d in results/*/; do
  # basename strips the path; wc -l counts the files listed
  echo "$(basename "$d"): $(ls "$d" | wc -l) files"
done
\end{lstlisting}

\section{Windows line endings}\label{windows-line-endings}

Windows ends each line of text with two invisible characters (carriage return, line feed), and Linux with one. A script written on Windows can fail on Linux with baffling errors such as \texttt{\$\textquotesingle{}\textbackslash{}r\textquotesingle{}:\ command\ not\ found}. Fix it on the cluster:

\begin{lstlisting}[language=bash]
# cluster
# In each line, replace a carriage return (\r) at the end of the line ($) with nothing
sed -i 's/\r$//' myscript.sh
\end{lstlisting}

